\chapter{Discussion}
\label{chap:discussion}

\section{Answer to the research question}

Within the evaluated boundary, independent runtime enforcement reduced
prompt-injection-driven violations while retaining useful task performance.
This is a bounded answer, not a claim that Sentinel solves prompt injection.
It rests on recurrence across two complementary settings rather than on any
single zero-attack result.

AgentDojo provides the broad controlled comparison: across all three matched
model configurations, enforcement lowered attacker success while raising both
utility and secure utility. WASP then tests whether the same pattern transfers
to a browser-native loop. Across its three seeds, full enforcement reduced
action-audited terminal success from 17/252 episodes to zero and
attacker-directed execution from 42/252 to zero. Legitimate completion changed
from 154/252 to 149/252, while secure utility rose from 138/252 to 149/252
(\cref{tab:agentdojo-e2e,tab:wasp-27b}). Between-seed variation and the fixed
benchmark scope prevent a stronger population claim, but the joint direction
is clear: the security change did not depend on making the agent inactive.

The two benchmarks reach that conclusion through different interfaces and
oracles, so their absolute rates are not pooled. Their agreement instead
supports the architectural thesis. A probabilistic planner can retain freedom
to choose useful actions without deciding alone what receives authority to
execute. Sentinel moves that authority to independently governed observation
and action boundaries, allowing model capability and execution permission to
remain separate questions.

\section{Interpreting security and utility}

The most consequential pattern is not simply that attacker success fell, but
that containment did not require inactivity. An indirect injection competes
with the user's task for context, steps, and future decisions. Removing hostile
influence can therefore recover useful work as well as suppress an attacker
goal. The AgentDojo gains are consistent with that explanation, whereas WASP's
small aggregate utility loss shows that the benefit is neither automatic nor
uniform. Because intervention timing was not isolated experimentally, this
remains an interpretation of the joint results rather than a separate causal
effect.

This distinction changes the appropriate counterfactual. Clean utility
estimates capability without an attack; it does not show what the same attacked
trajectory would have achieved without enforcement. That comparison comes from
the matched off/enforce or off/full pair. Secure utility then captures the joint
outcome in which the user's objective succeeds and the attacker's does not,
while paired gains and losses reveal regressions hidden by an aggregate rate.
Security and usefulness should therefore be reported together rather than
treated as opposing headline metrics.

The same joint view explains why terminal attack success is necessary but
insufficient. In the WASP baseline, substantially more episodes executed an
attacker-directed action than completed the final attacker objective. This gap
is the benchmark's ``security by incompetence'': the browser may already have
acted for the attacker even when the agent cannot finish the goal
~\cite{evtimov2025wasp}. Attempt and execution telemetry exposes those effects
and distinguishes a model proposal from an action that crossed the privileged
boundary. The replay remains conservative, however, because it cannot classify
reasoning-only diversion or every semantically equivalent action.

This interpretation also constrains comparisons across models. A more capable
agent may complete more legitimate and malicious goals, while a weak agent can
appear secure by failing at both. Since the evaluated configurations differ in
planning and recovery behaviour as well as size, neither parameter count nor a
low off-condition attack rate is an independent security axis. The defensible
comparison remains within each matched model and reports utility, terminal
outcomes, and executed effects together.

\section{Why both enforcement points matter}

The observation and action boundaries answer different questions. The first
changes what information reaches the planner; the second decides whether a
concrete effect may cross into the privileged executor. The policy-only/full
comparison makes that distinction visible. Policy-only retains boundary
neutralisation, a detector warning, and action policy, whereas full adds the
verified rewrite and suppresses the warning after a confirmed removal. At the
same aggregate utility of 149/252, episodes with attacker-directed proposals
and executions fell from 15/252 to zero; terminal attacker success changed only
from 1/252 to zero (\cref{tab:wasp-27b,tab:wasp-seeds}). The action audit
therefore reveals a trajectory change that the terminal oracle alone would
largely hide.

The historical 9B development sequence provides a useful negative result.
Retaining hostile content while warning the model did not improve terminal or
intermediate attack outcomes over off. Because middleware changed across that
sequence, it cannot identify the warning's independent effect. It nevertheless
shows why the final study treats warning around retained content as weaker
development evidence rather than as an equivalent observation boundary.

The action guard addresses the residual problem. A detector may miss an
unfamiliar instruction, a proposal may reflect earlier state, and an unrelated
model error can reach the same privileged tool. A policy decision immediately
before execution can still enforce stable invariants without reconstructing the
model's full intention. The few sensitive-data-flow blocks recorded after
AgentDojo sanitisation illustrate this backstop: observation governance was the
dominant intervention, but it did not make pre-tool control redundant.

Relying on either boundary alone would create a different weakness. Action-only
control would let hostile content consume context and planning steps, leaving
policy to recognise its effects one by one. Since legitimate and
attacker-directed work may use the same tools and origins, a broad policy risks
refusing useful actions while a narrow one risks missing disguised harmful
effects. Sanitisation alone avoids that trade-off for detected text but supplies
no independent authority over protected data, unsafe schemes, uploads, or
consequential actions. The two-boundary design instead places each decision
where its evidence is available: content provenance before planning, then
capability, target, and live state before execution
(\cref{sec:observation-governance,sec:action-governance}).

This composition also explains Sentinel's model independence. The planner need
not expose hidden reasoning or implement a security-specific API; it receives a
governed message and returns the ordinary BrowserUse tool proposal. Sentinel
can then evaluate deterministic context that would be unsuitable to place in
the prompt, including protected-value classes and local file inspection.
Model independence here means interface compatibility, not equal performance:
the model still determines how well the remaining task is planned.

\section{Policy, authority, and semantic limits}

The WASP tasks expose a basic difficulty for action policy: attacker and user
goals often converge on the same-origin forms and controls. Treating every
effect after a detector finding as approval-requiring would turn untrusted
provenance into broad denial of service. Sentinel instead applies deterministic
rules to independently sensitive sinks, including prohibited cross-origin
flow, protected values, unsafe schemes, local files, code, and configured
consequential targets. Post-verified sanitisation addresses the observation
that shaped the plan, while policy constrains the effect eventually proposed.

This division is practical rather than complete. A natural-looking instruction
can evade the detector and resemble a legitimate task step, while an allowed
same-origin click can still have consequences that are not visible in the
local DOM. Policy can therefore establish compliance with encoded constraints;
it cannot prove that every permitted effect advances an open-ended goal. Task
Shield, MELON, and DRIFT address this semantic problem through task-derived or
counterfactual abstractions
~\cite{jia2025taskshield,zhu2025melon,li2025drift}. A promising composition
would retain deterministic executor policy for non-negotiable invariants while
an independently isolated component checks ambiguous effects against a trusted
task representation.

Selective authority follows the same separation of concerns. A hard rule is
not an invitation to seek an override; a permissible but consequential
proposal may receive narrow, expiring approval; a protected operation is
performed directly by the operator; and an unresolved observation is
quarantined rather than approved. These states differ in both who acts and what
authority is granted. Server-side scope, live revalidation, and paired audit
keep that distinction intact even when the interface presents each state as an
interruption (\cref{sec:human-authority}).

Tests and walkthroughs establish that the implemented workflow fails closed,
revalidates stale proposals, and records its outcomes. They do not establish
that people will understand the distinctions, respond appropriately, or avoid
warning fatigue. The authority mechanism is therefore complete as a technical
boundary within the thesis scope, while its human effectiveness remains an
empirical question.

\section{Implications for system-level defences}

Seen against the systems in \cref{tab:system-comparison}, Sentinel occupies a
deliberate middle ground. It does not inherit the formal guarantees of CaMeL or
Fides because it preserves an existing free-form browser loop rather than
deriving execution from a trusted program. Its contribution is correspondingly
empirical: independently governed input rewriting, pre-tool enforcement, human
authority, and paired audit operated together while retaining useful task
performance. Structured planners offer stronger guarantees where their task
abstraction applies; an executor-side wrapper offers a practical boundary when
the agent exposes no such representation.

The comparison also suggests a useful hybrid. A structured task abstraction
could judge whether ambiguous content and effects remain relevant to the
user's goal, while Sentinel continues to enforce non-negotiable local rules and
bind human authority to live browser state. That arrangement could reduce
reliance on surface signatures without returning the final decision to the
acting model. Its value would depend on whether the semantic monitor adds
coverage without enough false intervention to erase the utility retained here.
This composition remains future work rather than an effect demonstrated by the
present experiments.

\section{Limitations and validity}
\label{sec:discussion-limitations}

Detector evidence has the clearest adaptation limitation. Several attack
families informed rule and representation changes during development, including
a WASP failure in which an unsafe link target was absent from BrowserUse's
accessibility text. The component scorecards therefore measure regression
coverage of known families, not generalisation to unseen attacks. The final
three-seed experiment has a stronger boundary because its middleware and
condition-specific policies were frozen before execution, but even a clean
post-check establishes only that the configured detector no longer fires. It
does not prove that the remaining text is free of hostile intent. A held-out
adaptive evaluation against the frozen whole system is the appropriate next
test; browser-specific work could follow MUZZLE by selecting injection surfaces
from agent behaviour and refining attacks through execution feedback
~\cite{zhan2025adaptive,syros2026muzzle}.

The evaluated model channel is textual. Screenshots were excluded from the
WASP observation path, so the results do not extend to pixel-only instructions
or tasks that require visual grounding. A visual extension would need its own
provenance, isolation, and action-authority design because images expand both
capability and attack surface
~\cite{bailey2024imagehijacks,koh2024visualwebarena}.

Experimental validity is also bounded by the selected models, benchmarks, and
execution design. AgentDojo contributes three model configurations, whereas
the final WASP comparison uses one 27B model, three sampling seeds, and a fixed
84-case matrix. Website state was reset before every cell, but the condition
order was not randomised; changes in service load or model-serving state could
therefore align with a later condition. The historical 9B WASP rows remain
development evidence, and some AgentDojo artefacts predate the final provenance
schema. In addition, AgentDojo enforcement combines external guards with fixed
system-level security guidance, so its rates describe the complete governed
configuration rather than the guards in isolation. These constraints preclude
deployment-population and causal model-scaling claims.

Finally, the policy and prototype boundaries are finite. Pattern-based
protected-data classification can miss a semantically sensitive value, while a
permitted same-origin action may trigger server-side consequences unavailable
to local inspection. Loopback origin checks are appropriate for the prototype
but do not isolate it from a compromised same-user process. Incidents, secret
takeover, live revalidation, extraction and history handling, and special-URL
controls have current test coverage, yet the frozen benchmark trajectories did
not measure their effects. They remain implementation assurance rather than
retrospective experimental evidence.

\subsection{Ethics and responsible research}
\label{subsec:ethics}

No human participated in the reported studies, and no participant data were
collected. Scripted benchmark modes supplied reproducible control behaviour, so
approval counts, pause duration, and Activity annotations are technical
telemetry. Interface walkthroughs establish server-owned request semantics and
rendered states, not comprehension, trust, workload, habituation, or recovery
quality.

The browser evidence came from controlled benchmark environments and local
synthetic interface fixtures. The submitted figures contain no live credential,
private account, private URL, or raw attack payload, and the dissertation does
not reproduce the benchmark's operational injection text. These choices reduce
unnecessary disclosure, although the work remains dual use because the analysis
can inform attacks as well as defences. Any future study of real operator
decisions would therefore require ethics approval before recruitment or data
collection, together with protocols for consent, data minimisation, and secure
handling of records.

\section{Implications and future work}

The limitations set priorities for further security evaluation. A frozen
adaptive study should test unseen prompt families and optimisation against the
combined boundaries on held-out tasks, using a larger seed set and randomised
condition order. A visual guard would require bounded, provenance-aware
representations and stricter authority for uncertain pixel-derived actions.
Meanwhile, an operator-visible task contract could state expected origins, data
classes, permitted effects, human steps, and stop conditions before exposure to
untrusted content. Each addition requires its own evaluation.

The authority workflow raises a separate human question. A participant study
should test whether operators understand interventions, distinguish among
approval, direct action, and incident review, and resume work without excessive
burden or warning fatigue. The implemented control plane is a technical basis
for such a study, not evidence of those outcomes. Recruitment would require
stable scenarios, validated measures, pilot testing, and the approved protocol
and safeguards described in \cref{subsec:ethics}.

Deployment beyond the local prototype would require authenticated control,
stronger process separation, protected policy and trace keys, and support for
transactional browser environments. This is distinct from benchmark robustness:
local isolation does not establish attack generalisation, just as a stronger
detector does not establish safe deployment.

The broader implication is therefore architectural. Under the tested
conditions, useful autonomy and independent authority were compatible:
Sentinel constrained what reached and left the planner without requiring a
perfect classifier or preventing the agent from completing useful work.
